\section{Introduction}

Predict-then-optimize pipelines are often evaluated by prediction error even
though the final objective is decision quality. Decision-focused learning (DFL)
addresses this mismatch by training predictors through an optimization-aware
surrogate such as SPO+ \citep{elmachtoub2022smart}. In shortest-path benchmarks,
this idea can produce large gains over a two-stage mean-squared-error baseline.
In our kidney-exchange experiments, however, the gains are more modest.

This paper asks a targeted question: when is DFL useful for graph combinatorial
optimization? Our starting point is that selected-solution identity is not
enough. A model can select a different feasible solution from the oracle and
still achieve nearly the same objective value. This can happen in decomposable
packing-style problems, where many feasible solutions are assembled from mostly
independent components and close substitutes are common.

We therefore analyze prediction errors at the decision level. Rather than asking
only which edge weights are predicted poorly, we ask whether the badly predicted
edges participate in the symmetric difference between selected and oracle
solutions, and whether the alternative solution is truly worse.

\paragraph{Contributions.}
We make three contributions. First, we validate the relevant SPO+ training path
against independent reference implementations so that later interpretation is
not confounded by implementation bugs. Second, we provide decision-level
diagnostics for kidney exchange that distinguish decision-critical prediction
errors from irrelevant prediction errors. Third, we build a toy family showing
that close-substitute behavior is not specific to one hand-picked instance: it
appears across several packing-style structures and contrasts with path-like
negative controls.
